\chapter{Conclusion}
\label{ch:conclusion}

\section{Problem}
Sugarcane farmers in Karnataka need region-specific recommendations in Kannada, and a
wrong pesticide or dose can destroy a crop or harm a person. A chatbot for them must
therefore be judged first on whether it ever gives a wrong recommendation, and only then
on how many questions it answers.

\section{Solution}
We built Krishi Mitra in two versions. Version~1 is a Kannada RAG pipeline with hybrid
BGE-M3 retrieval, an abstention gate, LLM generation and two answer checks. Version~2
answers only from typed fact records drafted from the UAS Bengaluru 2025 and UAS Dharwad
2020 Packages of Practices, each bound to its printed page, shown only after a person has
confirmed it against that page, and matched to the question without a language model.
Everything else is refused, with the Kisan Call Centre number.

\section{What was achieved}
\begin{itemize}
  \item A working Kannada chat assistant (Next.js interface, FastAPI backend) with two
    selectable answer paths.
  \item A set of safety failures found in a RAG pipeline and fixed, each with its effect:
    a confidence gate that read a rank-based score, a router label used as a hard filter,
    router and judge calls that failed silently when a reasoning model used up its token
    budget, and a numeric checker that bound numbers to the wrong units.
  \item Evidence that standard text metrics are unsuitable here: ROUGE-L scores 6 of 16
    Kannada references 0.000 against themselves, and chrF and BLEU score a 10$\times$
    overdose well above a correct paraphrase.
  \item A held-out evaluation of version~1 (29 of 40 correct decisions) and a no-retrieval
    baseline, which showed that retrieval reduces confident wrong answers but that a
    faithfulness judge cannot detect an answer about the wrong pest.
  \item A version~2 design in which that failure cannot occur by construction: an unknown
    pest name is refused, and no text is generated.
  \item 165 automated tests that protect these properties.
  \item Human review of 52 of the 90 records against the printed pages: 50 confirmed and
    usable, including early shoot borer, termite and weed-control pesticide records, and 2
    held because the registration check found carbendazim has no sugarcane entry in the
    CIB\&RC fungicide list.
  \item Two version~2 pilots with 30 new questions in total: 27 correct decisions, all 13
    should-refuse questions refused, and no wrong dose. The three errors were refusals: two
    from uncovered wording (since fixed) and one for a pest with no record.
  \item A verification label under every answer in place of an invented confidence score,
    and a prototype of the hybrid middle ground: an optional, tested query canonicaliser in
    which a language model can choose only a topic label and never reach a dose.
\end{itemize}

\section{Limitations}
\begin{itemize}
  \item \textbf{Coverage.} Version~2 can answer only what the two sugarcane chapters
    contain and only records a person has confirmed. 50 of 90 records were confirmed
    at the time of writing; the other 40, including 17 pest and disease records, are
    refused.
  \item \textbf{Review.} The fact records were checked by the project team, not yet by an
    agricultural expert. Pesticide records remain blocked until their CIB\&RC registration
    is confirmed.
  \item \textbf{Wording.} The matcher uses keyword lists. Questions phrased in ways the lists
    do not cover are refused (the safe direction), and symptom-based questions are refused
    by design. The Kannada keyword lists and messages need review by more Kannada speakers
    and by farmers.
  \item \textbf{Evaluation size.} The two version~2 pilots have 30 questions; with no wrong
    dose this bounds the wrong-dose rate only below about 10\%. The second pilot's writer had
    seen the records. The version~1 held-out questions were
    written with partial expert input and some AI assistance, and the development labels
    were assigned by the team, so all results are indicative.
  \item \textbf{Version~1 thresholds} were never calibrated; the development data showed that
    relevance alone cannot separate answerable from unanswerable questions.
  \item \textbf{Source age.} The UAS Dharwad book is from 2020; recommendations and pesticide
    registrations change over time.
\end{itemize}

\section{Future work}
\begin{itemize}
  \item \textbf{Complete confirmation of all 90 records.} 40 records are still unconfirmed,
    including 10 pest and 7 disease records; 23 of the 40 name a chemical product
    (12 chemicals in all), and carbendazim needs a registration source before its two
    sett-treatment records can be used. Each needs a person to check it against the printed page and to
    confirm its current CIB\&RC registration for sugarcane and that pest. The registration
    helper already lists, for each chemical, its sugarcane entries in the CIB\&RC lists and
    any banned or restricted-use entries; the entries found so far restrict use on other
    crops (for example chlorpyrifos on ber, citrus and tobacco), and the decision for
    sugarcane is the reviewer's.
  \item \textbf{A real-world evaluation set.} Replace the AI-assisted pilot questions with real
    farmer questions: Kisan Call Centre query logs (published on government open-data
    portals) filtered to Karnataka sugarcane, and transcriptions of real Kannada calls or of
    conversations recorded by village extension workers. Expected labels (answer or refuse,
    and the correct record) must be assigned by agricultural experts, not by the system or
    an AI tool, and the set must be large enough for a meaningful bound (for example 100
    questions with no failure bounds the failure rate below 3\%).
  \item \textbf{Hybrid middle grounds for understanding the question.} Keep the deterministic
    table lookup as the only source of doses, but use a model only to understand the
    question: (a) an LLM used solely for query canonicalisation or intent extraction, mapping
    a free-form Kannada question to one of the existing subject and topic identifiers in
    \texttt{intents.json}; and (b) constrained (grammar-based) decoding, so the model can only
    output a valid identifier or ``none'', never free text. The lookup, review status,
    registration check and refusal rules stay unchanged, so the model can cause a refusal or
    select the wrong topic but can never write a dose. Its key safety test is that an
    unknown pest must map to ``none'' and be refused, never to the nearest known pest (the
    version~1 failure). A first version of (a) and (b) is implemented, tested with a
    simulated model (Section~\ref{sec:canonicalizer}), and probed once with a live model on 19
    developer-written questions: 9 of 10 topics right and 9 of 9 must-refuse questions refused
    (Section~\ref{sec:canonprobe}). The next step is to evaluate it on real farmer questions
    the keyword lists miss, over repeated runs.
  \item Voice input and output in Kannada, and a WhatsApp channel, since many farmers use
    voice and messaging more than typed web chat.
  \item Re-evaluating the reranker and the RAG path on a larger, multi-crop collection, where
    first-stage retrieval is no longer near its ceiling.
\end{itemize}
